% \documentclass{cmspaper}
%\begin{document}

\section{Electron Studies} \label{sec:electrons}

\subsection{Identification and Isolation} \label{sec:IdAndIso}

An electron reconstructed object is formed from a deposit of energy in ECAL (called a ``supercluster'') with a track pointing towards it.
In this study the standard ``PixelMatchGsfElectrons'' collection is used \cite{GSFele}.
In this reconstruction algorithm the supercluster is first matched to hits in the pixel detector.
The pixel hits are then used to seed track finding.  Finally a track fit is performed using the Gaussian Sum Filter 
(GSF) algorithm, which allows the track to be reconstructed up to the ECAL surface
despite kinks due to radiated bremsstrahlung.  
The $P_{T}$ and $\eta$ of the electrons are calculated using both calorimeter and tracker informations, 
taking into account the different detector resolutions as a function of $P_{T}$.
Standard CMSSW corrections are applied to set correctly the electromagnetic (e.m.) energy scale.
%that is matched with hits in the pixel detector.  Few electrons loose momentum to radiated photons before the pixel detector since the majority of the material in the tracker is after the pixels.  Thus, the energy-weighted average impact point of the electrons and radiated photons with the supercluster will correspond well with the track found using these pixel hits as seeds.  

The criteria for electron identification (ID) at the reconstruction level are very loose, resulting in many fake electrons 
in the collection. In this analysis the HEEP selections for electron ID and isolation \cite{EleID}, which are optimized for 
electrons with an energy of hundreds of GeV, have been applied in an effort to 1) keep the efficiency high on electrons with high $P_{T}$ from LQ decays, and 
2) reduce the number of jets faking electrons in the collection.
The definitions of the variables are the following:
%
\begin{itemize}
%
\item $H/E$: the ratio of the hadronic energy of all the HCAL Rec Hits in a cone of $\Delta R=0.1$, centered on the 
electron's position in the calorimeter, to the electromagnetic energy of the supercluster associated to the electron;
%
\item $\sigma_{\eta\eta}$: the shape variable, which describes the electron shower size along the $\eta$ direction, is 
defined as 
\begin{displaymath}
\sigma_{\eta\eta} = \sum_{i \in \mathrm{crystals}} ( \eta_i - \eta_s )^2 \frac{E_i}{E_{\mbox{seed cluster}}} \quad ,
\end{displaymath}
where $i$ is the index of a crystal, and $s$ is the index of the seed cluster of the electron's supercluster;
%
\item $\Delta\eta_{in}$ and $\Delta\phi_{in}$: the difference in $\eta$ and $\phi$ between the track position as 
measured in the inner layer tracker, extrapolated to the ECAL surface, and the $\eta$ and $\phi$ of the supercluster;
%
\item number of tracks: with $P_{T}>1.5$ GeV in a cone of $0.1 < \Delta\mbox{R} < 0.2 $ around the electron direction;
%
\item track isolation: the sum of the $P_{T}$ of the tracks (as defined above) divided by the $E_{T}$ of the supercluster;
%
\item e.m. isolation: the transverse electromagnetic energy of all the clusters in a cone $\Delta\mbox{R} < 0.3$ centered on the 
electron's position in the calorimeter, excluding the clusters that make up the supercluster;
%
\item hadronic isolation: the transverse hadronic energy of all the HCAL Rec Hits with an $E_{T}>0$ in a cone of
$0.15 < \Delta\mbox{R} < 0.3$ centered on the electron's position in the calorimeter. 
%
\end{itemize}

Figure~\ref{fig:elecID} and Figure~\ref{fig:elecIso} show, respectively, the distributions of ID and isolation variables, 
for both true electrons coming from LQ decays, and fake electrons from the sample of QCD multi-jets events. 
Table~3
%\ref{tab:HEEPselection} 
summarizes the cuts from HEEP selection which are 
used in this analysis. In addition only reconstructed electrons with $P_{T}>30$ GeV and 
within the ECAL fiducial region 
$|\eta|<1.4$ + $1.6<|\eta|<2.0$ are included in the electron collection used in the analysis.  

Recentely the HEEP selection criteria have been slightly improved on e.m. and hadronic isolation, 
leading to a better rejection of jets faking electrons, and keeping the same efficiency on real electrons.
The impact of this improved selection criteria on the results here presented will be investigated in 
future upgrades of this analysis. 

%The quantities used in this definition are the ratio of hadronic to electromagnetic energy (H/E), shower shape along psudorapidity ($\sigma_{\eta\eta}$), difference in psudorapidity between the track and the supercluster ($\Delta\eta_{in}$), and the difference in phi between the track and supercluster ($\Delta\phi_{in}$). These cuts are appropriate for start up conditions as they allow a high efficiency.  The quatities and values used for these selection can be found in table~\ref{tab:ElecIdIso}. Figure~\ref{fig:elecID} shows the distributions of these variables for both true electrons coming from the leptoquark decays and fake electrons from a QCD sample.

%Further cleaning of the collection is achieved through requiring that the electron be isolated.  The three quantities considered are the number of tracks within a cone of radius 0.2 of the contact with the supercluster, the ratio of the sum of transverse energy of tracks within a cone of radius 0.2 to that of the electron track (Track Et ratio), and the ratio of transverse energy in the ECAL within a cone of radius 0.4 to that of the electron (ECAL Et ratio).  These are also found in table~\ref{tab:ElecIdIso}.  The distributions of these variables for both true electrons from the leptoquark decay and fake electrons from QCD events can be seen in figure~\ref{fig:elecIso}
 
%From : twiki.cern.ch/twiki/bin/view/CMS/SWGuideEgamma

% Robust Electron Def from : twiki.cern.ch/twiki/bin/view/CMS/SWGuideElectronID
% Ele Isolation from : twiki.cern.ch/twiki/bin/view/CMS/SWGuideEgamma 

%Still need description of these plots, like what cuts were used...

%\begin{2figures}{hbtp}
%  \resizebox{\linewidth}{!}{\includegraphics{plots/h_eleDeltaEtaTrkSC.eps}} &
%  \resizebox{\linewidth}{!}{\includegraphics{plots/h_eleDeltaPhiTrkSC.eps}} \\
%  \caption{eleDeltaEtaTrkSC}
%  \label{fig:eleDeltaEtaTrkSC} &
%  \caption{\small \sl eleDeltaPhiTrkSC}
%  \label{fig:eleDeltaPhiTrkSC} \\
%\end{2figures}
%\begin{2figures}{hbtp}
%  \resizebox{\linewidth}{!}{\includegraphics{plots/h_eleHoE.eps}} &
%  \resizebox{\linewidth}{!}{\includegraphics{plots/h_eleSigmaEE.eps}} \\
%  \caption{\small \sl eleHoE}
%  \label{fig:eleHoE} &
%  \caption{\small \sl eleSigmaEE}
%  \label{fig:eleSigmaEE} \\
%\end{2figures}
%\begin{2figures}{hbtp}
%  \resizebox{\linewidth}{!}{\includegraphics{plots/h_eleNumTrkIso.eps}} &
%  \resizebox{\linewidth}{!}{\includegraphics{plots/h_eleEcalIso.eps}} \\
%  \caption{\small \sl eleNumTrkIso}
%  \label{fig:eleNumTrkIso} &
%  \caption{\small \sl eleEcalIso}
%  \label{fig:eleEcalIso} \\
%\end{2figures}
%\begin{figure}
%  \begin{center}
%    \resizebox{0.5\linewidth}{!}{\includegraphics{plots/h_eleTrkIso.eps}}
%    \caption{\small \sl eleTrkIso}
%    \label{fig:eleTrkIso}
%  \end{center}
%\end{figure}

\begin{figure}
  \begin{center}
  \begin{tabular}{cc}
  \resizebox{7cm}{!}{\includegraphics{plots/h_eleHoE.eps}} &
  \resizebox{7cm}{!}{\includegraphics{plots/h_eleSigmaEE.eps}} \\
  \resizebox{7cm}{!}{\includegraphics{plots/h_eleDeltaEtaTrkSC.eps}} &
  \resizebox{7cm}{!}{\includegraphics{plots/h_eleDeltaPhiTrkSC.eps}} \\
  \end{tabular} 
  \caption{\small \sl Distribution of electron ID variables for true electrons coming from the decay of a 650 GeV mass LQ, and for fake electrons from 
the sample of QCD multijets events. Reconstructed electrons with $P_{T}>150$ 
GeV are showed. Reconstructed electrons from LQ decay are matched with generated electrons if $\Delta$R between them is less than 0.02. 
The left peak in the $\sigma_{\eta\eta}$ distribution corresponds to electrons falling in the 
ECAL barrel, the right peak to those falling in the endcap.}
 \label{fig:elecID}
  \end{center}
\end{figure}

\begin{figure}
  \begin{center}
  \begin{tabular}{cc}
  \resizebox{7cm}{!}{\includegraphics{plots/h_eleNumTrkIso.eps}} &
  \resizebox{7cm}{!}{\includegraphics{plots/h_eleTrkIsoAbs.eps}}\\
  \resizebox{7cm}{!}{\includegraphics{plots/h_eleEcalIsoAbs.eps}} &
  \resizebox{7cm}{!}{\includegraphics{plots/h_eleHcalIsoAbs.eps}} \\
  \end{tabular}
  \caption{\small \sl Distribution of electron isolation variables for true electrons coming from the decay of a 650 GeV mass LQ, 
    and for fake electrons from the sample of QCD multijets events. Reconstructed electrons with $P_{T}>150$ GeV are showed. 
    Reconstructed electrons from LQ decay are matched with generated electrons if $\Delta$R between them is less than 0.02.}
    \label{fig:elecIso}
  \end{center}
\end{figure}

\begin{table}[htbp]
  \label{tab:HEEPselection}
  \begin{center}
    \begin{tabular}{|lcc|lcc|} \hline
      \multicolumn{3}{|c|}{ID Variables} & \multicolumn{3}{|c|}{Isolation Variables} \\ 
      Variable & Barrel & Endcap & Variable & Barrel & Endcap  \\ \hline
      $H/E$  & $<0.05$ & $<0.08$ & number of tracks  & $<4$ & $<4$ \\ \hline
      $\sigma_{\eta\eta}$  & $<0.011$ & $<0.0275$ & track isolation (GeV) & $<7.5$ & $<15$ \\ \hline
      $|\Delta\eta_{in}|$  & $<0.005$ & $<0.007$ & e.m. isolation (GeV) & $<6+0.01*E_{t}$ & $<6+0.01*E_{t}$ \\ \hline
      $|\Delta\phi_{in}|$  & $<0.09$ & $<0.09$ & hadronic isolation (GeV) & $<4+0.005*E_{t}$ & $<4+0.005*E_{t}$ \\ \hline
    \end{tabular}
  \caption{\small \sl HEEP electron identification and isolation selection criteria}
  \end{center}
\end{table}


\subsection{Acceptance and Reconstruction Efficiency}

Electrons coming from leptoquark decays, 
being typically isolated from jets and having a $P_{T}$ of hundreds of GeV can be easily reconstructed with high efficiency. 
Table~4
%\ref{tab:ElecEffAcc} 
shows the electron acceptance in the fiducial region ($A_{FR}$), the 
reconstruction efficiency within the fiducial region ($\varepsilon_{FR}^{reco}$), and the overall reconstruction efficiency 
($A_{FR} \times \varepsilon_{FR}^{reco}$) for single electrons coming from LQ decays for different LQ masses.
The overall reconstruction efficiency on a single electron, including the acceptance cuts, is around $70\%$ for LQ masses between 250 GeV and 1 TeV.
As an illustration, Figure \ref{fig:elecEffFV} shows the reconstruction efficiency of an electron
from a LQ of mass 650 GeV as a function of $\eta$. The plot has been done using the FastSim 
due to the higher statistics available in FastSim samples with respect to FullSim ones. 
%Inside the fiducial region, the reconstruction efficiency is high, and only slightly dependent on $\eta$. 
A relevant loss of efficiency is seen in the region around $|\eta| \approx 1.5$ due to the detector crack between the ECAL barrel and the endcaps, 
which is excluded from the analyis. %INEFFICIENCY AT HIGH ETA? 

\begin{table}[htb]
  \label{tab:ElecEffAcc}
  \begin{center}
    \begin{tabular}{|l|c|c|c|} \hline
      LQ mass (GeV) & $A_{FR}$ & $\varepsilon_{FR}^{reco}$ & $A_{FR} \times \varepsilon_{FR}^{reco}$\\ \hline
      250 & 80.8\% & 88.1\% & 71.2\% \\ \hline
      400 & 84.4\% & 87.1\% & 73.5\% \\ \hline
      650 & 87.3\% & 86.6\% & 75.6\% \\ \hline
      1000 & 89.2\% & 83.1\% & 74.1\% \\ \hline
    \end{tabular}
    \caption{\small \sl Electron acceptance in the fiducial region ($A_{FR}$), the 
      reconstruction efficiency within the fiducial region ($\varepsilon_{FR}^{reco}$), and the overall reconstruction efficiency 
      ($A_{FR} \times \varepsilon_{FR}^{reco}$) for single electrons coming from LQ decays. Efficiencis are 
      calculated using FullSim samples for different LQ masses. Statistical relative errors on efficiency are less than 2\%.}
  \end{center}
\end{table}

\begin{figure}
  \begin{center}
    \resizebox{9cm}{!}{\includegraphics{plots/Ele_eff_M650GeV.eps}}
    \caption{\small \sl Electron reconstruction efficiency as a function of $\eta$. The filled area represents the  
      ECAL regions out of the fiducial volume defined at \ref{sec:IdAndIso}. The plot is done using a 
      LQ sample of 650 GeV mass produced with FastSim.}
    \label{fig:elecEffFV}
  \end{center}
\end{figure}

% indico.cern.ch/getFile.py/access?contribId=3&resId=0&materialId=slides&confId=32048

\subsection{Comparison of Electron Variables between FastSim and FullSim} \label{sec:elecFastSimCompare}
%\subsection{Kinematic Distributions for Signal Electrons}
An effort is made to compare reconstructed quantities between FastSim and FullSim for high energy electrons.
The comparison is performed using LQ samples produced with FastSim in CMSSW\_1\_6\_9. 
Several kinematics variables, ID and isolation variables, and efficiencies have been investigated. 
Figure \ref{fig:elecVariables} shows the distributions of some of these quantities for reconstructed electrons from LQ decay 
for both FastSim and FullSim.
A generally good agreement within a few \% is found in the shape of all the distributions studied for reconstructed electrons matched 
with real electrons from LQ decay. Some discrepancies of up to 5\% are observed in the electron 
reconstruction efficiency,  which is overestimated in FastSim. 
These studies will be repeated in future with more recent versions of FastSim in order to verify if the level 
of agreement with FullSim, in particular for the efficiency, has been improved. 

\begin{figure}
  \begin{center}
  \begin{tabular}{c}
    %\resizebox{9cm}{!}{\includegraphics{plots/Ele_eff_M650GeV.eps}}
    \resizebox{10cm}{!}{\includegraphics{plots/pT_eta_FastVsFull_M650.eps}} \\
    \resizebox{10cm}{!}{\includegraphics{plots/DPhi_DEta_FastVsFull_M650.eps}} \\ 
%  \resizebox{5cm}{!}{\includegraphics{plots/UMD.eps}} &
%  \resizebox{5cm}{!}{\includegraphics{plots/UMD.eps}} \\
  \end{tabular}
  \caption{\small \sl Distributions of some reconstructed quantities for electrons from LQ decay, 
    obtained with FastSim (blue filled histogram) and FullSim (black dots) 
    samples with $M_{LQ}=650$ GeV.
    From the top-left to the bottom-right: $P_{T}$, $\eta$, $\Delta\phi_{in}$, $\Delta\eta_{in}$.}
  \label{fig:elecVariables}
  \end{center}
\end{figure}

%efficiency (pT, eta) (no deltaR from genJet since already showed distance OK above)
%pT, eta dist
%{\Large \sl Energy Resolution}
%fast vs full

%\end{document}
